\begin{figure}[ht]
\centering
\begin{tikzpicture}[x=1cm, y=1cm, font=\small]

% Axes
\draw[->, thick] (-0.3, 0) -- (8.2, 0) node[right] {$x$};
\draw[->, thick] (0, -0.3) -- (0, 5.5) node[above] {$y = f(x)$};

% Curve f(x) = 0.18 (x-1)^2 + 0.6
\draw[thick, engblue] plot[domain=0.3:7.6, samples=80]
  (\x, {0.18*(\x-1)*(\x-1) + 0.6});

% Point of tangency a = 3
\fill[engred] (3, {0.18*4 + 0.6}) circle (2pt);
\node[engred, anchor=north east] at (3, {0.18*4 + 0.6}) {$(a,\,f(a))$};

% Secant 1 (long h)
\draw[thick, engamber, dashed] (1.0, {0.18*0 + 0.6}) -- (7.4, {0.18*40.96 + 0.6});
\fill[engamber] (7.0, {0.18*36 + 0.6}) circle (1.6pt);
\node[engamber, anchor=south west, font=\scriptsize] at (7.0, {0.18*36 + 0.6 + 0.05}) {$h_{1}$ large};

% Secant 2 (medium h)
\draw[thick, enggray, dashed] (2.0, {0.18*1 + 0.6}) -- (6.0, {0.18*25 + 0.6});
\fill[enggray] (5.5, {0.18*20.25 + 0.6}) circle (1.4pt);
\node[enggray, anchor=south west, font=\scriptsize] at (5.5, {0.18*20.25 + 0.6 + 0.05}) {$h_{2}$};

% Tangent line at a = 3: slope = 0.36*(a-1) = 0.72
\draw[thick, engred] (0.5, {0.78 - 0.72*2.5}) -- (6.0, {0.78 + 0.72*3.0});
\node[engred, anchor=west] at (5.7, {0.78 + 0.72*2.7 + 0.2}) {tangent at $a$};

% Bracket showing h shrinking
\draw[<->, thick] (3, -0.5) -- (5.5, -0.5);
\node[anchor=north, font=\scriptsize] at (4.25, -0.5) {$h \to 0$: secants rotate into tangent};

\end{tikzpicture}
\caption[Tangent line as limit of secants]{The difference quotient
$(f(a + h) - f(a))/h$ is the slope of the secant line through
$(a, f(a))$ and $(a + h, f(a + h))$. As $h$ shrinks, the second
point slides along the graph toward $(a, f(a))$ and the secant
rotates about the fixed point. The derivative $f'(a)$ is the slope
of the limiting line, the tangent at $a$.}
\label{fig:vol02:ch05:tangent-secant}
\end{figure}
